\documentclass{article}

\usepackage{arxiv}
\usepackage[utf8]{inputenc}
\usepackage[T1]{fontenc}
\usepackage{hyperref}
\usepackage{url}
\usepackage{booktabs}
\usepackage{amsfonts}
\usepackage{nicefrac}
\usepackage{microtype}
\usepackage{graphicx}
\usepackage{multirow}
\usepackage{array}
\usepackage{xcolor}
\usepackage{float}

\title{Benchmarking PyCaret AutoML Against BiLSTM for Fine-Grained Emotion Classification: A Comparative Study on 20-Class Emotion Detection}

\author{
  Arya Muda Siregar \\
  Institut Teknologi Sumatera \\
  Lampung Selatan, Indonesia \\
  \texttt{123450063@student.itera.ac.id} \\
  \And
  Arielva Simon Siahaan \\
  Institut Teknologi Sumatera \\
  Lampung Selatan, Indonesia \\
  \texttt{123450105@student.itera.ac.id} \\
  \And
  Haikal Fransisko Simbolon \\
  Institut Teknologi Sumatera \\
  Lampung Selatan, Indonesia \\
  \texttt{123450106@student.itera.ac.id} \\
  \AND
  Martin C.T. Manullang\thanks{Corresponding author} \\
  Institut Teknologi Sumatera \\
  Lampung Selatan, Indonesia \\
  \texttt{martin.manullang@itera.ac.id} \\
}

\begin{document}
\maketitle

\begin{abstract}
Fine-grained emotion classification—identifying specific emotional states such as happiness, anger, sadness, and fear—is a challenging yet important task in Natural Language Processing (NLP). This paper presents a benchmark comparison between classical Machine Learning (ML) approaches and a Deep Learning (DL) approach for multi-class emotion classification over 20 emotion categories. On the ML side, we evaluate Logistic Regression, Multinomial Naive Bayes, and Support Vector Machine (LinearSVC) using TF-IDF feature representations. On the DL side, we compare Bidirectional Long Short-Term Memory (BiLSTM), Gated Recurrent Unit (GRU), and a lightweight Transformer model implemented in PyTorch. Experiments are conducted on the 20-Emotion Text Classification Dataset comprising 79,595 English sentences. Results show that the BiLSTM model achieves the best overall performance with an accuracy of 89\% and a weighted F1-score of 0.89, outperforming both the best ML model (SVM: 88.11\%) and other DL variants. Our findings demonstrate that while traditional ML models remain competitive and computationally efficient, sequence-based deep learning models better capture contextual emotional cues in text.
\end{abstract}

\keywords{Emotion Classification \and Sentiment Analysis \and BiLSTM \and Support Vector Machine \and PyCaret \and Natural Language Processing \and Deep Learning}

\newpage
\section{Introduction}

Emotion recognition in text is a fundamental challenge in NLP with broad applications ranging from mental health monitoring and customer feedback analysis to social media sentiment tracking. Unlike binary sentiment analysis (positive/negative), fine-grained emotion classification requires distinguishing between a larger and more nuanced set of emotional states, significantly increasing the complexity of the task \cite{ekman1992emotions}.

Recent progress in NLP has produced two dominant paradigms for tackling text classification: traditional Machine Learning (ML) methods relying on handcrafted features, and Deep Learning (DL) approaches that learn representations directly from raw text. Classical ML methods such as Support Vector Machines and Logistic Regression have demonstrated strong performance on many text classification benchmarks, especially when paired with effective feature extraction techniques such as TF-IDF \cite{joachims1998text}. Meanwhile, deep learning architectures like Long Short-Term Memory networks \cite{hochreiter1997long}, Gated Recurrent Units \cite{cho2014gru}, and Transformer models \cite{vaswani2017attention} have shown remarkable ability to model sequential dependencies in text.

Despite extensive research, there is limited direct comparison between these two paradigms on fine-grained emotion classification tasks with large class sets (20 classes). This paper aims to bridge this gap by conducting a systematic benchmark study that evaluates both ML and DL approaches on the same dataset under identical experimental conditions.

The contributions of this work are as follows:
\begin{itemize}
    \item A systematic benchmark of three ML classifiers (Logistic Regression, Naive Bayes, SVM) using TF-IDF features for 20-class emotion classification.
    \item A comparative evaluation of three DL architectures (BiLSTM, GRU, Transformer) implemented in PyTorch on the same dataset.
    \item An analysis of the trade-offs between model performance, training time, and computational complexity across all six models.
    \item Deployment of the best ML and DL models as interactive web applications on Hugging Face Spaces.
\end{itemize}

\section{Related Work}

Early approaches to emotion recognition relied heavily on lexicon-based methods and handcrafted features \cite{strapparava2008semeval}. These methods, while interpretable, often fail to generalize across domains. The introduction of machine learning methods, particularly Support Vector Machines, provided significant improvements in classification accuracy by learning discriminative boundaries in high-dimensional feature spaces \cite{pang2002thumbs}.

The advent of deep learning marked a paradigm shift in NLP. Hochreiter and Schmidhuber \cite{hochreiter1997long} proposed LSTM networks capable of capturing long-range dependencies in sequential data—a critical property for understanding the contextual nuances of emotion in text. The Bidirectional LSTM \cite{schuster1997bidirectional} further improved upon this by processing sequences in both forward and backward directions, allowing the model to incorporate both past and future context simultaneously.

Cho et al. \cite{cho2014gru} introduced the Gated Recurrent Unit as a simplified variant of LSTM, reducing computational overhead while retaining strong sequential modeling capability. The Transformer architecture \cite{vaswani2017attention}, based entirely on self-attention mechanisms, has since become the dominant paradigm in NLP, enabling parallelization and scalability.

For emotion classification specifically, Saravia et al. \cite{saravia2018carer} demonstrated the effectiveness of LSTM-based models on Twitter emotion data. Mohammad and Turney \cite{mohammad2013crowdsourcing} contributed foundational work on crowdsourcing emotion annotations for NLP benchmarks. More recently, AutoML frameworks such as PyCaret \cite{ali2020pycaret} have democratized access to ML pipelines by automating model selection, hyperparameter tuning, and evaluation, making it easier to compare multiple classifiers on the same task.

\section{Dataset}

\subsection{Description}

We use the \textbf{20-Emotion Text Classification Dataset}, an English-language benchmark comprising 79,595 sentences labeled with one of 20 distinct emotion categories. The emotion labels include: \textit{happiness, sadness, anger, fear, love, surprise, disgust, trust, anticipation, anxiety, boredom, confusion, curiosity, excitement, guilt, loneliness, nostalgia, pride, relief,} and \textit{shame}.

\subsection{Statistics}

\begin{table}[H]
\centering
\caption{Dataset Statistics}
\begin{tabular}{ll}
\toprule
\textbf{Property} & \textbf{Value} \\
\midrule
Total Samples & 79,595 \\
Number of Classes & 20 \\
Language & English \\
Task & Multi-class Text Classification \\
Average Sentence Length & \textasciitilde{}12 tokens \\
\bottomrule
\end{tabular}
\end{table}

\subsection{Preprocessing}

Prior to feature extraction and model training, we apply the following text preprocessing steps:
\begin{enumerate}
    \item \textbf{Lowercasing}: All text is converted to lowercase to reduce vocabulary size.
    \item \textbf{Punctuation and special character removal}: Non-alphanumeric characters are removed.
    \item \textbf{Tokenization}: Text is split into individual tokens.
    \item \textbf{Stopword removal}: Common English stopwords are removed for ML feature extraction.
    \item \textbf{Train/Validation/Test Split}: The dataset is split into 80\% training, 10\% validation, and 10\% test sets with stratified sampling to maintain class distribution.
\end{enumerate}

For ML models, TF-IDF (Term Frequency-Inverse Document Frequency) vectors are computed with a maximum vocabulary size of 50,000 terms. For DL models, we use a word embedding layer with an embedding dimension of 128 and a vocabulary size of 30,000.

\section{Methodology}

\subsection{Machine Learning Pipeline}

We employ TF-IDF feature extraction to transform raw text into numerical feature vectors, which are then fed into three ML classifiers. We adopt the AutoML benchmark approach to compare multiple models on the same feature representation:

\textbf{Logistic Regression (LR)}: A linear classifier optimized using the L-BFGS solver with L2 regularization. Used as the baseline model due to its simplicity and interpretability.

\textbf{Multinomial Naive Bayes (MNB)}: A probabilistic classifier based on Bayes' theorem, well-suited for text classification with TF-IDF features. Applied with Laplace smoothing ($\alpha = 1.0$).

\textbf{Support Vector Machine (LinearSVC)}: A linear SVM optimized using the squared hinge loss with L2 regularization. Known for strong performance on high-dimensional text classification tasks \cite{joachims1998text}.

\subsection{Deep Learning Architecture}

All DL models are implemented in PyTorch \cite{paszke2019pytorch} with the following shared configuration:

\begin{itemize}
    \item Embedding dimension: 128
    \item Vocabulary size: 30,000
    \item Maximum sequence length: 128 tokens
    \item Dropout rate: 0.3
    \item Optimizer: Adam with learning rate 0.001
    \item Batch size: 64
    \item Training epochs: 10
\end{itemize}

\textbf{BiLSTM (Bidirectional Long Short-Term Memory)}: Two stacked LSTM layers (hidden size = 256) processing input sequences in both forward and backward directions. The final hidden states from both directions are concatenated and passed through a fully connected layer with softmax activation.

\textbf{GRU (Gated Recurrent Unit)}: Two stacked GRU layers (hidden size = 256) with a bidirectional configuration. The GRU architecture reduces the number of gate operations compared to LSTM while maintaining comparable performance.

\textbf{Transformer}: A lightweight encoder-only Transformer with 4 attention heads, 2 encoder layers, and a feed-forward dimension of 512. A positional encoding layer is added to inject sequential information.

All DL models contain fewer than 10 million parameters, conforming to the parameter budget constraint of this study.

\section{Experiments}

\subsection{Experimental Configuration}

All experiments are conducted on Google Colab with GPU acceleration (NVIDIA Tesla T4). ML models are trained on CPU using scikit-learn \cite{pedregosa2011scikit}, while DL models utilize GPU training via PyTorch \cite{paszke2019pytorch}.

\subsection{Evaluation Metrics}

We report the following metrics for all models:
\begin{itemize}
    \item \textbf{Accuracy}: Proportion of correctly classified samples.
    \item \textbf{Macro F1-Score}: Unweighted mean F1-score across all 20 classes.
    \item \textbf{Weighted F1-Score}: F1-score weighted by class support.
    \item \textbf{Training Time}: Wall-clock time for model training.
\end{itemize}

\subsection{Hyperparameter Settings}

\begin{table}[H]
\centering
\caption{Hyperparameter Configuration for DL Models}
\begin{tabular}{lcc}
\toprule
\textbf{Hyperparameter} & \textbf{BiLSTM / GRU} & \textbf{Transformer} \\
\midrule
Hidden Size & 256 & 128 \\
Layers & 2 & 2 \\
Attention Heads & -- & 4 \\
Dropout & 0.3 & 0.3 \\
Optimizer & Adam & Adam \\
Learning Rate & 0.001 & 0.001 \\
Batch Size & 64 & 64 \\
Epochs & 10 & 10 \\
\bottomrule
\end{tabular}
\end{table}

\section{Results and Discussion}

\subsection{Benchmark Results}

Table~\ref{tab:benchmark} presents the benchmark results for all six models evaluated on the test set.

\begin{table}[H]
\centering
\caption{Benchmark Results: ML vs. DL Models on 20-Emotion Classification}
\label{tab:benchmark}
\begin{tabular}{llccc}
\toprule
\textbf{Type} & \textbf{Model} & \textbf{Accuracy} & \textbf{F1-Score} & \textbf{Training Time} \\
\midrule
\multirow{3}{*}{ML} 
  & Logistic Regression & 0.8656 & -- & -- \\
  & Naive Bayes         & 0.8381 & -- & -- \\
  & \textbf{SVM (LinearSVC)} & \textbf{0.8811} & \textbf{0.8808} & -- \\
\midrule
\multirow{3}{*}{DL}
  & \textbf{BiLSTM}     & \textbf{0.89} & \textbf{0.89} & 5m 53s \\
  & GRU                 & 0.88 & 0.88 & 7m 09s \\
  & Transformer         & 0.87 & 0.86 & 11m 09s \\
\bottomrule
\end{tabular}
\end{table}

\subsection{Discussion}

\textbf{ML Models}: Among the three ML classifiers, SVM (LinearSVC) achieves the best performance with an accuracy of 88.11\% and a weighted F1-score of 0.8808. This is consistent with findings in the literature that SVMs excel at text classification in high-dimensional TF-IDF feature spaces \cite{joachims1998text}. Logistic Regression provides competitive performance (86.56\%), while Naive Bayes lags behind at 83.81\%, likely due to its strong feature independence assumption being violated in emotion-laden text.

\textbf{DL Models}: BiLSTM achieves the highest overall performance (accuracy: 89\%, weighted F1: 0.89) among all models. The bidirectional processing allows the model to capture both preceding and following context, which is particularly beneficial for detecting subtle emotional cues that depend on word order and sentence structure \cite{schuster1997bidirectional}. GRU (88\% accuracy) achieves comparable results to BiLSTM while requiring fewer parameters. The Transformer model (87\% accuracy) underperforms relative to recurrent models in this setting, likely due to the relatively moderate dataset size which may not be sufficient to fully leverage attention-based representations without pre-training \cite{vaswani2017attention}.

\textbf{ML vs. DL Comparison}: The BiLSTM DL model outperforms the best ML model (SVM) by approximately 0.9 percentage points in accuracy, demonstrating the advantage of learned sequential representations over bag-of-words TF-IDF features. However, the performance gap is relatively small, suggesting that classical ML methods with strong feature engineering remain highly competitive for this task. Notably, the SVM model requires significantly less training infrastructure and produces results comparable to DL models without the need for GPU acceleration.

\textbf{Training Efficiency}: ML models converge almost instantaneously compared to DL models. BiLSTM trains in approximately 6 minutes, GRU in 7 minutes, and the Transformer in 11 minutes—all on GPU. For deployment scenarios where compute resources are limited, SVM offers a favorable trade-off between accuracy and efficiency.

\section{Conclusion}

This paper presents a systematic benchmark study comparing Machine Learning and Deep Learning approaches for fine-grained emotion classification across 20 emotion categories. Our experiments on the 20-Emotion Text Classification Dataset (79,595 samples) demonstrate that:

\begin{itemize}
    \item \textbf{BiLSTM is the best overall model}, achieving the highest accuracy (89\%) and weighted F1-score (0.89) among all six evaluated models.
    \item \textbf{Deep Learning outperforms Machine Learning} in both accuracy and


\bibliographystyle{plain}
\bibliography{reference}

\end{document}
